\section{Ablation Study}
\label{sec:ablation}

We isolate the contribution of each \method\ component on a representative subset
of the standard (largely stationary) benchmarks---ETTh1, ETTm2, and
Weather---at horizons $H \in \{96, 720\}$ with lookback $L{=}336$ and three
seeds. Because every component is an opt-in generalization that degenerates to a
known baseline (Section~\ref{sec:method}), each ablation is a clean controlled
comparison along the nesting $\text{Linear} \subseteq \text{RLinear} \subseteq
\method$.

\subsection{Component Ablations}
\label{sec:ablation:components}

Table~\ref{tab:ablation} reports the ablations. We consider: (A0) the full
\method; (A1) replacing \arevin\ with plain RevIN while keeping the decomposition,
which isolates the value of adaptive de-normalization; (A2) replacing the
learnable spectral split with a DLinear-style fixed moving-average split while
keeping \arevin, which isolates the value of the learnable decomposition; (A3)
$K{=}1$ with \arevin\ (an ``A-RevIN-Linear''), showing the decomposition's
marginal value; (A4) $K{=}1$ with plain RevIN, which is exactly RLinear and
verifies the degeneracy $\method \supseteq \text{RLinear}$; (A5) increasing the
number of bands to $K \in \{3,4\}$; (A7) freezing the decomposition at
initialization to separate the benefit of a soft mask from that of
\emph{learning} it; and (A8) \arevin\ without the drift-propagation term
$\boldsymbol{\lambda}$.

\IfFileExists{../results/tables/ablation_table.tex}{%
  % Auto-generated by scripts/make_tables.py.
\begin{table}[t]
\centering
\caption{Ablation study (mean test MSE/MAE over the subset \{ETTh1, ETTm2, Weather\}$\times H\in\{96,720\}\times$3 seeds, $L=336$). Each variant degenerates a single FreqLite component.}
\label{tab:ablation}
\begin{tabular}{lrrr}
\toprule
Variant & Avg MSE & Avg MAE & Avg Params \\
\midrule
A0: Full FreqLite & 0.310 & 0.342 & 276,221 \\
A1: $-$A-RevIN (plain RevIN) & 0.310 & 0.342 & 274,996 \\
A2: fixed MA split & 0.310 & 0.343 & 276,219 \\
A3: $K{=}1$ + A-RevIN & 0.312 & 0.345 & 138,723 \\
A4: $K{=}1$ + RevIN ($\equiv$RLinear) & 0.312 & 0.345 & 137,498 \\
A5: $K{=}3$ bands & 0.309 & 0.342 & 413,719 \\
A5: $K{=}4$ bands & \textbf{0.309} & 0.342 & 551,217 \\
A6: gated recombination & 0.310 & 0.343 & 276,223 \\
A7: frozen decomposition & 0.310 & 0.342 & 276,219 \\
A8: $-\lambda$ drift term & 0.310 & 0.342 & 275,813 \\
\bottomrule
\end{tabular}
\end{table}
%
}{\begin{table}[t]\centering\caption{Component ablations.}%
  \label{tab:ablation}\textit{[ablation table pending]}\end{table}}

\subsection{Analysis}
\label{sec:ablation:analysis}

The A4 row provides the key sanity check: with a single all-pass band and plain
RevIN, \method\ reproduces RLinear (mean MSE $0.3119$), confirming the strict
nesting that makes every added component an independent, opt-in generalization.

On this stationary subset, the small overall improvement of the full model
(A0, $0.3096$) over RLinear (A4, $0.3119$)---about $0.7\%$---is attributable to
the \emph{band split}, not to adaptive normalization. Two comparisons make this
precise. Adding \arevin\ to a single-head backbone leaves the error essentially
unchanged (A3 $0.3118$ vs.\ A4 $0.3119$), and replacing \arevin\ with plain RevIN
in the full model is likewise neutral (A1 $0.3096$ vs.\ A0 $0.3096$); so
\arevin\ is neutral on stationary data---exactly the graceful-degradation
behavior it is designed for. By contrast, introducing the band split accounts for
the gain: a fixed moving-average split is nearly as good (A2 $0.3102$), and
freezing the learned split at initialization is no worse than learning it here
(A7 $0.3096$), indicating that on these datasets it is the \emph{soft band
routing} rather than the \emph{learning} of the cutoff that matters. Adding more
bands gives only a marginal further change (A5 $0.3094$ for $K\in\{3,4\}$), and
removing the drift term has no effect on this stationary subset (A8 $0.3096$).
These results motivate the dedicated non-stationarity study in
Section~\ref{sec:results:nonstat}, where---unlike here---\arevin\ is the
component that drives the gains.

\subsection{Interpretability of Learned Parameters}
\label{sec:ablation:interp}

Because \method's added components are few and scalar-parameterized, their
learned values are directly interpretable. The learned spectral mask
(Eq.~\ref{eq:mask-low}) is shown in Figure~\ref{fig:learned_filter}, and the
per-step \arevin\ correction profile $\lambda_t$ (Eq.~\ref{eq:arevin-shift}) in
Figure~\ref{fig:arevin_profile}. Together with the gate analysis of
Section~\ref{sec:results:nonstat}---where the learned gate $\bar\rho$ opens on
non-stationary data and the error falls monotonically with gate
engagement---these provide direct evidence that the components behave as designed
rather than merely adding parameters.
